% ============================================================
%  appendix_examples.tex -- A catalogue of worked vertices
% ============================================================

\chapter{A Catalogue of Worked Vertices}
\label{app:examples}

This appendix collects a set of vertices derived with FeynPy, ordered
roughly by increasing complexity, as declaration/output pairs. Every
expression below is the verbatim printed output of \py{feynman_rule(...)}
with \py{simplify=True}; nothing has been reformatted beyond inserting line
breaks. The purpose is twofold: to give a reader an idea of what the output
actually looks like before committing to \Cref{ch:feynpy}, and to make the
notation of \Cref{sec:validation-canonical} concrete.

\section*{Reading the output}

Four printed heads account for almost everything.

\begin{center}
\begin{tabular}{@{}P{4.3cm}p{7.2cm}@{}}
\toprule
\normalfont\textbf{Printed head} & \textbf{Meaning} \\
\midrule
\texttt{mink(4, mu1)}   & Lorentz index of leg $1$ \\
\texttt{bis(4, i1)}     & spinor (bispinor) index of leg $1$ \\
\texttt{cof(3, c1)}, \texttt{coad(8, a1)} & colour fundamental and adjoint indices \\
\texttt{g(...)}         & metric, or a Kronecker delta on a gauge index \\
\texttt{t(...)}, \texttt{f(...)} & $\mathrm{SU}(N)$ generator and structure constant \\
\texttt{gamma(...)}, \texttt{gamma5(...)} & $\gamma^\mu$ and $\gamma^5$ \\
\texttt{pcomp(q1, mu2)} & component $\mu_2$ of the incoming momentum $q_1$ \\
\bottomrule
\end{tabular}
\end{center}

\noindent Legs are numbered in the order they are passed to
\py{feynman_rule}, all momenta are incoming, and every rule carries the
overall factor $+i$ of \Cref{tab:theory-conventions}. Index names ending in
\texttt{\_int} or \texttt{\_canon} are generated dummies and carry no
meaning.

% --------------------------------------------------------------
\section{Scalar \texorpdfstring{$\phi^4$}{phi\string^4}}
% --------------------------------------------------------------

\begin{minted}{python}
Phi = Field("Phi", spin=0, self_conjugate=True, symbol=S("phi"))
L   = Model(lam * Phi * Phi * Phi * Phi).lagrangian()
L.feynman_rule(Phi, Phi, Phi, Phi)
\end{minted}

\begin{minted}{text}
24I*lam
\end{minted}

\noindent The factor $24 = 4!$ is the number of ways of assigning four
identical fields to four legs, produced by the permutation sum of
\Cref{sec:internals-contraction} rather than inserted by hand.

% --------------------------------------------------------------
\section{QED}
% --------------------------------------------------------------

\begin{minted}{python}
Psi    = Field("Psi", spin=Fraction(1,2), self_conjugate=False,
               symbol=S("psi"), conjugate_symbol=S("psibar"),
               indices=(SPINOR_INDEX,), quantum_numbers={"Q": S("Qpsi")})
Photon = Field("A", spin=1, self_conjugate=True,
               symbol=S("A"), indices=(LORENTZ_INDEX,))
U1QED  = GaugeGroup(name="U1QED", abelian=True, coupling=S("e"),
                    gauge_boson=Photon, charge="Q")

L = Model(I * Psi.bar * Gamma(mu) * DC(Psi, mu),
          gauge_groups=(U1QED,)).lagrangian()
L.feynman_rule(Psi.bar, Psi, Photon)
\end{minted}

\begin{minted}{text}
I*Qpsi*e*gamma(bis(4, i1),bis(4, i2),mink(4, mu3))
\end{minted}

\noindent That is $i Q_\psi e\, \gamma^{\mu_3}$, the vertex derived by hand
in \Cref{sec:feynrules-algorithm}. Note that the interaction was never
written: it came out of the expansion of \py{DC}.

% --------------------------------------------------------------
\section{Scalar QED}
% --------------------------------------------------------------

\begin{minted}{python}
# Photon and U1QED are those of the QED example above
PhiC = Field("PhiC", spin=0, self_conjugate=False,
             symbol=S("phiC"), conjugate_symbol=S("phiCdag"),
             quantum_numbers={"Q": S("Qphi")})

L = Model(DC(PhiC.bar, mu) * DC(PhiC, mu),
          gauge_groups=(U1QED,)).lagrangian()

L.feynman_rule(PhiC.bar, PhiC, Photon)            # current
L.feynman_rule(PhiC.bar, PhiC, Photon, Photon)    # seagull
\end{minted}

\begin{minted}{text}
-I*e*Qphi*pcomp(q1,mu3) + I*e*Qphi*pcomp(q2,mu3)

2I*e^2*Qphi^2*g(mink(4, mu3),mink(4, mu4))
\end{minted}

\noindent The first is $-i e Q_\phi (q_1 - q_2)^{\mu_3}$; the second is the
seagull $2 i e^2 Q_\phi^2 g^{\mu_3\mu_4}$, which arises from the term
quadratic in the gauge field produced when both covariant derivatives are
expanded. One declared term, two vertices of different arity.

% --------------------------------------------------------------
\section{QCD}
% --------------------------------------------------------------

The declaration is the one given in full in \Cref{sec:feynpy-qcd-example}.
Its four terms produce the following.

\subsection*{Quark--gluon}

\begin{minted}{text}
I*gs*gamma(bis(4, i1),bis(4, i2),mink(4, mu3))
   *t(coad(8, a3),cof(3, c1),cof(3, c2))
\end{minted}

\noindent $i g_s \gamma^{\mu_3} t^{a_3}_{c_1 c_2}$.

\subsection*{Three-gluon}

\begin{minted}{text}
-gs*g(mink(4,mu3),mink(4,mu1))*f(coad(8,a1),coad(8,a2),coad(8,a3))
    *pcomp(q1,mu2)
+gs*g(mink(4,mu3),mink(4,mu1))*f(coad(8,a1),coad(8,a2),coad(8,a3))
    *pcomp(q3,mu2)
+gs*g(mink(4,mu3),mink(4,mu2))*f(coad(8,a1),coad(8,a2),coad(8,a3))
    *pcomp(q2,mu1)
-gs*g(mink(4,mu3),mink(4,mu2))*f(coad(8,a1),coad(8,a2),coad(8,a3))
    *pcomp(q3,mu1)
+gs*g(mink(4,mu1),mink(4,mu2))*f(coad(8,a1),coad(8,a2),coad(8,a3))
    *pcomp(q1,mu3)
-gs*g(mink(4,mu1),mink(4,mu2))*f(coad(8,a1),coad(8,a2),coad(8,a3))
    *pcomp(q2,mu3)
\end{minted}

\noindent Collecting the six terms in pairs gives the textbook form
\[
    g_s f^{a_1a_2a_3}\Big[
      g^{\mu_1\mu_2}(q_1-q_2)^{\mu_3}
    + g^{\mu_2\mu_3}(q_2-q_3)^{\mu_1}
    + g^{\mu_3\mu_1}(q_3-q_1)^{\mu_2}\Big].
\]

\subsection*{Four-gluon}

\begin{minted}{text}
-I*gs^2*g(mink(4,mu3),mink(4,mu4))*g(mink(4,mu1),mink(4,mu2))
       *f(coad(8,a),coad(8,a1),coad(8,a3))*f(coad(8,a),coad(8,a2),coad(8,a4))
-I*gs^2*g(mink(4,mu3),mink(4,mu4))*g(mink(4,mu1),mink(4,mu2))
       *f(coad(8,a),coad(8,a2),coad(8,a3))*f(coad(8,a),coad(8,a1),coad(8,a4))
-I*gs^2*g(mink(4,mu3),mink(4,mu1))*g(mink(4,mu4),mink(4,mu2))
       *f(coad(8,a),coad(8,a3),coad(8,a4))*f(coad(8,a),coad(8,a1),coad(8,a2))
+I*gs^2*g(mink(4,mu3),mink(4,mu1))*g(mink(4,mu4),mink(4,mu2))
       *f(coad(8,a),coad(8,a2),coad(8,a3))*f(coad(8,a),coad(8,a1),coad(8,a4))
+I*gs^2*g(mink(4,mu3),mink(4,mu2))*g(mink(4,mu4),mink(4,mu1))
       *f(coad(8,a),coad(8,a3),coad(8,a4))*f(coad(8,a),coad(8,a1),coad(8,a2))
+I*gs^2*g(mink(4,mu3),mink(4,mu2))*g(mink(4,mu4),mink(4,mu1))
       *f(coad(8,a),coad(8,a1),coad(8,a3))*f(coad(8,a),coad(8,a2),coad(8,a4))
\end{minted}

\noindent The summed adjoint index \texttt{a} is the contraction between the
two structure constants. This vertex is a good illustration of why the
canonicalisation of \Cref{sec:validation-canonical} needs a controlled use of
the Jacobi identity: the same expression can be written with different pairs
of $f$'s, and only a fixed convention makes two derivations comparable.

\subsection*{Ghost--gluon}

\begin{minted}{text}
gs*f(coad(8, a1),coad(8, a3),coad(8, a2))*pcomp(q1,mu3)
\end{minted}

\noindent $g_s f^{a_1 a_3 a_2} q_1^{\mu_3}$, carrying the momentum of the
antighost leg, as expected from $\mathcal{L}_{\mathrm{gh}}$ in
\Cref{eq:theory-ghost}.

\subsection*{The gluon two-point function}

\begin{minted}{text}
 I*g(coad(8,a1),coad(8,a2))*pcomp(q1,mu1)*pcomp(q2,mu2)/xi
+I*g(mink(4,mu1),mink(4,mu2))*g(coad(8,a1),coad(8,a2))
    *pcomp(q1,mu1_int)*pcomp(q2,mu1_int)
-I*g(coad(8,a1),coad(8,a2))*pcomp(q1,mu2)*pcomp(q2,mu1)
\end{minted}

\noindent The gauge parameter $\xi$ appears explicitly, so setting $\xi = 1$
recovers the Feynman-gauge inverse propagator.

% --------------------------------------------------------------
\section{Flavour-expanded Yukawa}
% --------------------------------------------------------------

\begin{minted}{python}
Gen = flavor_index("Generation", 3, prefix="f")
lep = Field("l", spin=Fraction(1,2), self_conjugate=False,
            indices=(Gen, SPINOR_INDEX), symbol=S("l"),
            conjugate_symbol=S("lbar"), flavor_index=Gen,
            class_members=("e", "mu", "ta"))
Y   = Parameter("Y", indices=(Gen, Gen))
e, mu_, ta = lep.class_members

L = Model(Y(f1, f2) * lep.bar(f1) * lep(f2) * Phi).lagrangian()
\end{minted}

\noindent Without expansion the rule is generation-generic,

\begin{minted}{text}
I*g(bis(4, i1),bis(4, i2))*Y(f1,f2)
\end{minted}

\noindent and with \py{flavor_expand=True} the same declaration yields nine
explicit vertices:

\begin{minted}{text}
ebar  e   Phi -> I*g(bis(4,i1),bis(4,i2))*Y(1,1)
ebar  mu  Phi -> I*g(bis(4,i1),bis(4,i2))*Y(1,2)
ebar  ta  Phi -> I*g(bis(4,i1),bis(4,i2))*Y(1,3)
mubar e   Phi -> I*g(bis(4,i1),bis(4,i2))*Y(2,1)
mubar mu  Phi -> I*g(bis(4,i1),bis(4,i2))*Y(2,2)
mubar ta  Phi -> I*g(bis(4,i1),bis(4,i2))*Y(2,3)
tabar e   Phi -> I*g(bis(4,i1),bis(4,i2))*Y(3,1)
tabar mu  Phi -> I*g(bis(4,i1),bis(4,i2))*Y(3,2)
tabar ta  Phi -> I*g(bis(4,i1),bis(4,i2))*Y(3,3)
\end{minted}

\noindent Had \texttt{Y} been declared with \py{components} setting its
off-diagonal entries to zero, the six off-diagonal rules would simply not
exist.

% --------------------------------------------------------------
\section{Electroweak vertices from the packaged Standard Model}
% --------------------------------------------------------------

The vertices below come from \py{build_standard_model()}
(\Cref{sec:feynpy-models}), so they are already in the physical basis
produced by the field transformations of
\Cref{sec:feynpy-transformations}.

\begin{minted}{python}
sm = build_standard_model()
L  = sm.lagrangian
F  = sm.fields
\end{minted}

\subsection*{Triple gauge coupling $A W^- W^+$}

\begin{minted}{text}
 I*ee*g(mink(4,mu1),mink(4,mu2))*pcomp(q1,mu3)
-I*ee*g(mink(4,mu1),mink(4,mu2))*pcomp(q2,mu3)
-I*ee*g(mink(4,mu1),mink(4,mu3))*pcomp(q1,mu2)
+I*ee*g(mink(4,mu1),mink(4,mu3))*pcomp(q3,mu2)
+I*ee*g(mink(4,mu2),mink(4,mu3))*pcomp(q2,mu1)
-I*ee*g(mink(4,mu2),mink(4,mu3))*pcomp(q3,mu1)
\end{minted}

\noindent Structurally identical to the three-gluon vertex above, with the
structure constant replaced by the electric charge $e$ --- which is what the
$W^3$--$B$ mixing of \Cref{eq:theory-mixing} does to the
$\mathrm{SU}(2)_L$ self-coupling.

\subsection*{Higgs--$W$--$W$ and the triple Higgs coupling}

\begin{minted}{text}
L.feynman_rule(F.H, F.W.bar, F.W)  ->  I/2*vev*ee^2*g(mink(4,mu2),mink(4,mu3))/sw^2

L.feynman_rule(F.H, F.H, F.H)      ->  -6I*lam*vev
\end{minted}

\noindent Both are proportional to the vacuum expectation value, as they must
be: they exist only because $\Phi$ was shifted by $v$ in
\Cref{eq:theory-vev}.

\subsection*{Charged current $W^+ \bar u d$}

\begin{minted}{text}
 I/4*(1/2)^(1/2)*ee*g(cof(3,c2),cof(3,c3))
    *gamma(bis(4,i2),bis(4,i3),mink(4,mu1))*CKM(fl2,fl3)/sw
-I/4*(1/2)^(1/2)*ee*g(cof(3,c2),cof(3,c3))
    *gamma(bis(4,i2),bis(4,bis_canon_2),mink(4,mu1))
    *gamma5(bis(4,bis_canon_2),bis(4,i3))*CKM(fl2,fl3)/sw
+ ... (two further terms)
\end{minted}

\noindent The CKM matrix appears exactly where
\Cref{sec:theory-ssb} says it should, contracted on the two quark flavour
indices. The four terms are the expansion of the chiral projector pair
$P_L \otimes P_L$ into $\tfrac14(1 \mp \gamma^5)\otimes(1\mp\gamma^5)$, which
is the printed form the toolkit uses; recombining them is the chirality
question examined in \Cref{sec:validation-chirality}. The neutral currents
$A \bar\ell \ell$ and $Z \bar\nu \nu$ have the same shape, with the
electric-charge and $s_W/c_W$ combinations in place of the CKM matrix.
